\begin{frame}{``정의를 먼저''의 역사: 세 시대}
  \small
  ``먼저 문제를 정의하라''가 추상적 교훈이 아니라 실제 역사에서 어떤
  결과를 만들었는지 세 시대로:
  \begin{itemize}
    \item \kb{1959 사무엘의 체커스 (Samuel, 1959)}: ``머신러닝''이라는
      말을 현대적 의미로 만든 논문. 기계를 ``플레이하게'' 하는 게
      아니라, ``\textbf{이 판국이 얼마나 좋은가}(평가함수, $y$)를
      \textbf{학습}하는 문제''로 \textbf{재정의} --- 판국 수가 사람이
      규칙을 전부 쓸 수 없을 만큼 많으므로 기계가 self-play에서 배우게.
      (이 재정의가 40년 뒤 \textbf{AlphaGo 계보}의 출발점.)
    \item \kb{1970s 전문가 시스템 (MYCIN, 1976) --- 역의 길}:
      의사들이 ``X 증상이면 Y를 의심하라''는 \textbf{규칙을 수백 개
      직접} 써넣은 감염 진단. 좁은 영역에서 전문의에 필적했으나,
      규칙 \textbf{유지보수가 사람에게}, 새 영역마다 수천 개를 다시
      사람이 --- 교훈: ``정의(=모든 규칙을 사람이 쓰는 것)가 사람에
      의존할수록, 상한도 사람의 역량에 갇히고 유지비도 비싸진다.''
    \item \kb{1998 LeNet과 우편 자동화 (Week 10)}: ``입력: 손글씨
      숫자 이미지, 출력: 어느 숫자인가''. 사람이 형상(입력/출력)과
      특징을 설계하고, ``특징을 \textbf{어떻게 쓸 것인가}''는
      \textbf{기계가 데이터에서 학습} --- USPS 우편 분류에서
      ``사람은 문제만 정의하고 \textbf{기계가 세부 규칙을 배운다}''
      경로가 \textbf{확장성에 이긴 첫 대규모 증거}.
  \end{itemize}
\end{frame}
